% ============================================================================
% LANGENTWURF: Verbos en -ar (Regelmäßige -ar Verben)
% ============================================================================
% Sequenz 4: Mi colegio | Block c: Verbos en -ar
% Fachfarbe: #c41e3a (Karminrot)
% ============================================================================

\documentclass[11pt,a4paper]{article}

% ============================================================================
% PACKAGES
% ============================================================================
\usepackage[utf8]{inputenc}
\usepackage[T1]{fontenc}
\usepackage[ngerman]{babel}
\usepackage{geometry}
\usepackage{fancyhdr}
\usepackage{lastpage}
\usepackage{xcolor}
\usepackage{tcolorbox}
\usepackage{enumitem}
\usepackage{tabularx}
\usepackage{longtable}
\usepackage{array}
\usepackage{booktabs}
\usepackage{hyperref}
\usepackage{ifthen}
\usepackage{amssymb}

% ============================================================================
% KONFIGURATION
% ============================================================================

\newcommand{\plantier}{1}
\newcommand{\fachid}{spanisch}
\newcommand{\fach}{Spanisch}
\newcommand{\fachfarbe}{c41e3a}

% ============================================================================
% VARIABLEN
% ============================================================================

\newcommand{\jahrgangsstufe}{7}
\newcommand{\schulart}{Gesamtschule}
\newcommand{\stundenthema}{Verbos en -ar}
\newcommand{\reihenthema}{Mi colegio}
\newcommand{\stundennummer}{4c}
\newcommand{\reihenstunden}{6}
\newcommand{\unterrichtsdatum}{}
\newcommand{\unterrichtszeit}{}
\newcommand{\raum}{}

\newcommand{\lehrkraft}{N.N.}
\newcommand{\ausbildungsstand}{Lehrkraft}

\newcommand{\lerngruppengroesse}{24}
\newcommand{\maennlich}{12}
\newcommand{\weiblich}{12}
\newcommand{\divers}{0}

\newcommand{\gerniveau}{A1}
\newcommand{\lehrwerk}{Qué pasa 1}
\newcommand{\lehrwerkseite}{S. 48--51}

% ============================================================================
% LAYOUT
% ============================================================================
\geometry{left=2.5cm, right=2.5cm, top=2.5cm, bottom=2.5cm}
\definecolor{fach}{HTML}{\fachfarbe}
\definecolor{gray}{HTML}{666666}
\definecolor{lightgray}{HTML}{f5f5f5}
\definecolor{successgreen}{HTML}{2a7a4b}
\definecolor{warningorange}{HTML}{b35c00}

\tcbuselibrary{skins,breakable}

\newtcolorbox{infobox}[1][]{
    colback=lightgray,
    colframe=fach,
    fonttitle=\bfseries,
    title=#1,
    breakable
}

\newtcolorbox{scorebox}{
    colback=white,
    colframe=fach,
    fonttitle=\bfseries,
    title=Didaktik-Score,
    breakable
}

\pagestyle{fancy}
\fancyhf{}
\fancyhead[L]{\small\textcolor{gray}{\fach{} -- Jg. \jahrgangsstufe{} -- \stundenthema}}
\fancyhead[R]{\small\textcolor{gray}{Langentwurf Tier 1}}
\fancyfoot[C]{\small\thepage{} / \pageref{LastPage}}
\renewcommand{\headrulewidth}{0.4pt}
\renewcommand{\footrulewidth}{0pt}

% ============================================================================
% DOKUMENT
% ============================================================================
\begin{document}

% ============================================================================
% DECKBLATT
% ============================================================================
\begin{titlepage}
    \centering
    \vspace*{2cm}

    {\large\textcolor{gray}{\schulart}}\\[2cm]

    {\Huge\textcolor{fach}{\textbf{Langentwurf}}}\\[0.5cm]
    {\large Tier 1 -- Vollständig}\\[2cm]

    {\LARGE\textbf{\stundenthema}}\\[0.5cm]
    {\large Regelmäßige Verben auf -ar}\\[0.5cm]
    {\large im Rahmen der Unterrichtsreihe}\\[0.3cm]
    {\Large\textit{\reihenthema}}\\[2cm]

    \begin{tabular}{rl}
        \textbf{Fach:} & \fach{} (\gerniveau) \\[0.3cm]
        \textbf{Klasse:} & \jahrgangsstufe \\[0.3cm]
        \textbf{Lehrwerk:} & \lehrwerk, \lehrwerkseite \\[0.3cm]
        \textbf{Datum:} & \underline{\hspace{3cm}} \\[0.3cm]
        \textbf{Zeit:} & \underline{\hspace{3cm}} (Doppelstunde) \\[0.3cm]
        \textbf{Raum:} & \underline{\hspace{3cm}} \\[0.3cm]
    \end{tabular}

    \vfill

    {\large Lehrkraft: \lehrkraft}\\[0.3cm]
    {\normalsize\textcolor{gray}{\ausbildungsstand}}\\[1cm]

\end{titlepage}

% ============================================================================
% INHALTSVERZEICHNIS
% ============================================================================
\tableofcontents
\newpage

% ============================================================================
% 1. BEDINGUNGSANALYSE
% ============================================================================
\section{Bedingungsanalyse}

\subsection{Lerngruppe}
\begin{tabular}{ll}
    Kursgröße: & \lerngruppengroesse{} SuS (\maennlich{} m / \weiblich{} w / \divers{} d) \\
    Jahrgangsstufe: & \jahrgangsstufe \\
    Sprachniveau: & \gerniveau{} (GER) \\
    Fremdsprache: & 2. Fremdsprache \\
\end{tabular}

\subsection{Differenzierung (intern)}

\textbf{WICHTIG:} Keine Sternchen auf Schülermaterial!

\begin{tabular}{|l|p{4cm}|p{4cm}|p{4cm}|}
\hline
& \textbf{[B] Basis} & \textbf{[S] Standard} & \textbf{[E] Erweiterung} \\
\hline
\textbf{Ziel} & Berufsreife & Mittlere Reife & Gymnasium \\
\hline
\textbf{Inhalt} & yo, tú, él/ella & + nosotros & + vosotros, ellos \\
\hline
\textbf{Hilfen} & Konjugationstabelle, Beispiele & Teils vorgegeben & Nur Impulse \\
\hline
\end{tabular}

\subsection{Besonderheiten}
\begin{itemize}
    \item \textbf{Vorkenntnisse:} SuS kennen bereits \textit{ser} (unregelmäßig)
    \item \textbf{LRS:} 2 SuS (Zeitzugabe, farbige Endungen)
    \item \textbf{DaZ:} 1 SuS (deutsche Verbkonjugation parallel)
    \item \textbf{Leistungsstark:} 4 SuS (Zusatzaufgaben, alle Personen)
\end{itemize}

% ============================================================================
% 2. SACHANALYSE
% ============================================================================
\section{Sachanalyse}

\subsection{Sprachliche Strukturen}

\textbf{Konjugation regelmäßiger -ar Verben:}

\begin{center}
\begin{tabular}{|l|c|c|l|}
\hline
\textbf{Person} & \textbf{Endung} & \textbf{hablar} & \textbf{Deutsch} \\
\hline
yo & -o & hablo & ich spreche \\
tú & -as & hablas & du sprichst \\
él/ella/usted & -a & habla & er/sie spricht \\
nosotros/-as & -amos & hablamos & wir sprechen \\
vosotros/-as & -áis & habláis & ihr sprecht \\
ellos/ellas/ustedes & -an & hablan & sie sprechen \\
\hline
\end{tabular}
\end{center}

\textbf{Zielverben dieser Stunde:}
\begin{itemize}
    \item \textbf{hablar} -- sprechen, reden
    \item \textbf{estudiar} -- lernen, studieren
    \item \textbf{escuchar} -- hören, zuhören
    \item \textbf{trabajar} -- arbeiten
    \item \textbf{cantar} -- singen
    \item \textbf{bailar} -- tanzen
\end{itemize}

\subsection{Kontrastive Betrachtung (Deutsch $\leftrightarrow$ Spanisch)}
\begin{itemize}
    \item \textbf{Positive Transfers:} Ähnliche Konjugationslogik (Person-Endung), Wortstellung SVO
    \item \textbf{Interferenzen:} Subjektpronomen im Spanischen optional (,,Hablo español'' ohne ,,Yo'')
    \item \textbf{Falsche Freunde:} Keine bei den Zielverben
\end{itemize}

\subsection{Kulturelle Aspekte}
\begin{itemize}
    \item Schulalltag in Spanien (horario escolar)
    \item Typische Aktivitäten im Unterricht und in der Freizeit
\end{itemize}

% ============================================================================
% 3. DIDAKTISCHE ANALYSE
% ============================================================================
\section{Didaktische Analyse}

\subsection{Stellung in der Sequenz}
Dies ist die \textbf{\stundennummer. Doppelstunde} der Sequenz ,,\reihenthema'' (\reihenstunden{} Doppelstunden gesamt).

\textbf{Vorherige Stunden:}
\begin{itemize}
    \item 4a: Mi colegio -- Schulvokabular (Räume, Fächer)
    \item 4b: El horario -- Stundenplan, Uhrzeiten
\end{itemize}

\textbf{Folgestunden:}
\begin{itemize}
    \item 4d: En clase -- Vertiefung, Dialoge
    \item 4e: Projekt: Mi día escolar
\end{itemize}

\subsection{Kompetenzziele (Rahmenplan MV)}
\begin{itemize}
    \item \textbf{Kommunikativ:} Über den Schulalltag sprechen, einfache Aussagen machen
    \item \textbf{Sprachlich:} Regelmäßige -ar Konjugation (yo, tú, él/ella, nosotros)
    \item \textbf{Interkulturell:} Schulalltag in spanischsprachigen Ländern
    \item \textbf{Methodisch:} Konjugationsmuster erkennen und anwenden
\end{itemize}

\subsection{Legitimation}
\textbf{Rahmenplan Spanisch MV (Kl. 7-10):}
\begin{itemize}
    \item Kompetenzbereich: Verfügen über sprachliche Mittel -- Grammatik
    \item Standard: SuS können Verben im Präsens konjugieren
    \item Themenfeld: Schule und Lernen
\end{itemize}

% ============================================================================
% 4. STUNDENZIELE
% ============================================================================
\section{Stundenziele}

\subsection{Grobziel}
Die SuS erweitern ihre \textbf{grammatische Kompetenz}, indem sie regelmäßige -ar Verben konjugieren und in einfachen Sätzen über den Schulalltag anwenden.

\subsection{Feinziele (operationalisiert)}
\begin{tabular}{|c|p{8cm}|c|c|}
\hline
\textbf{Nr.} & \textbf{Die SuS können...} & \textbf{AFB} & \textbf{Diff.} \\
\hline
FZ1 & die Konjugationsendungen -o, -as, -a, -amos nennen & I & alle \\
\hline
FZ2 & die Verben hablar, estudiar, escuchar, trabajar, cantar, bailar konjugieren & II & alle \\
\hline
FZ3 & Sätze über den Schulalltag mit -ar Verben bilden & II & [S]/[E] \\
\hline
FZ4 & einen kurzen Text über ihren Schulalltag schreiben & III & [E] \\
\hline
\end{tabular}

% ============================================================================
% 5. METHODISCHE ANALYSE
% ============================================================================
\section{Methodische Analyse}

\subsection{Phasierung (Doppelstunde = 2 × 45 Min.)}
\begin{enumerate}
    \item \textbf{1. Stunde:}
    \begin{itemize}
        \item Einstieg: Bildimpuls (Schüler bei verschiedenen Aktivitäten)
        \item Erarbeitung: Induktive Grammatikeinführung (hablar als Modell)
        \item Übung: Chorisches Sprechen, Konjugationstabelle ergänzen
    \end{itemize}
    \item \textbf{2. Stunde:}
    \begin{itemize}
        \item Wiederholung: Schnelles Abfragen der Endungen
        \item Anwendung: Sätze bilden mit den sechs Zielverben
        \item Sicherung: Arbeitsblatt, Partner-Interview
    \end{itemize}
\end{enumerate}

\subsection{Medien}
\begin{itemize}
    \item Tafel/Whiteboard (Konjugationstabelle)
    \item Lehrwerk: \lehrwerk, \lehrwerkseite
    \item Arbeitsblatt: AB-04c.pdf
    \item Lernpfad: LP-04c.html (optional, digital)
    \item Flashcards: Verbinfinitive mit Bildern
\end{itemize}

% ============================================================================
% 6. VERLAUFSPLANUNG
% ============================================================================
\section{Verlaufsplanung}

\textbf{1. Stunde (45 Min.)}

\begin{longtable}{|p{1cm}|p{1.5cm}|p{4cm}|p{3cm}|p{2cm}|p{2cm}|}
\hline
\textbf{Zeit} & \textbf{Phase} & \textbf{Lehrerhandeln} & \textbf{Schülerhandeln} & \textbf{SF} & \textbf{Medien} \\
\hline
\endhead

0--5 & Einstieg & Begrüßung, zeigt Bilder von Aktivitäten & Beschreiben, vermuten & Plenum & Bilder \\
\hline

5--15 & Input & Präsentiert ,,hablar'', schreibt Konjugation an Tafel, erklärt Muster & Hören, notieren & Plenum & Tafel \\
\hline

15--25 & Übung 1 & Chorisches Sprechen anleiten: ,,hablo, hablas, habla...'' & Nachsprechen, rhythmisch & Plenum & -- \\
\hline

25--35 & Übung 2 & Verteilt AB, leitet Aufgabe 1 an & Endungen einsetzen & EA & AB \\
\hline

35--40 & Sicherung & Ergebnisse abfragen, Fehler korrigieren & Vergleichen & Plenum & Tafel \\
\hline

40--45 & Über\-leitung & Einführung weiterer Verben (estudiar, escuchar) & Notieren & Plenum & Tafel \\
\hline
\end{longtable}

\vspace{1em}

\textbf{2. Stunde (45 Min.)}

\begin{longtable}{|p{1cm}|p{1.5cm}|p{4cm}|p{3cm}|p{2cm}|p{2cm}|}
\hline
\textbf{Zeit} & \textbf{Phase} & \textbf{Lehrerhandeln} & \textbf{Schülerhandeln} & \textbf{SF} & \textbf{Medien} \\
\hline
\endhead

0--5 & Wieder\-holung & Schnelles Endungen-Quiz mündlich & Antworten rufen & Plenum & -- \\
\hline

5--15 & Input & Einführung trabajar, cantar, bailar & Notieren, nachsprechen & Plenum & Tafel \\
\hline

15--30 & Anwen\-dung & Leitet AB Aufg. 2-3 an, unterstützt & Verben konjugieren, Sätze bilden & EA/PA & AB \\
\hline

30--40 & Festigung & Partner-Interview: ,,¿Qué estudias?'' & Fragen und antworten & PA & -- \\
\hline

40--45 & Abschluss & Zusammenfassung, HA erklären & Notieren, aufräumen & Plenum & Tafel \\
\hline
\end{longtable}

\textbf{Legende:} EA = Einzelarbeit, PA = Partnerarbeit, SF = Sozialform

% ============================================================================
% 7. ERWARTETE SCHWIERIGKEITEN
% ============================================================================
\section{Erwartete Schwierigkeiten}

\begin{tabular}{|p{5cm}|p{8cm}|}
\hline
\textbf{Schwierigkeit} & \textbf{Reaktion} \\
\hline
Verwechslung -as / -a & Farbliche Markierung, Personalpronomen betonen \\
\hline
Vergessen des Stamms & Stamm unterstreichen, Formel ,,Infinitiv minus -ar'' \\
\hline
Aussprache hablar (stummes h) & Chorisches Sprechen, Einzelkorrektur \\
\hline
Nosotros-Form vergessen & Merkspruch: ,,Wir enden auf -amos'' \\
\hline
Heterogenität & Differenzierte AB-Aufgaben ([B]/[S]/[E]) \\
\hline
\end{tabular}

% ============================================================================
% 8. HAUSAUFGABE
% ============================================================================
\section{Hausaufgabe}

\begin{itemize}
    \item AB-04c Aufgabe 4 fertigstellen (Schulalltag beschreiben)
    \item Vokabeln lernen: hablar, estudiar, escuchar, trabajar, cantar, bailar
    \item Optional: Lernpfad LP-04c.html bearbeiten
\end{itemize}

% ============================================================================
% 9. ANLAGEN
% ============================================================================
\section{Anlagen}

\begin{enumerate}
    \item Arbeitsblatt (AB-04c.pdf)
    \item Musterlösung (ML-04c.pdf)
    \item Lehrerhinweise (LH-04c.pdf)
    \item Lernpfad (LP-04c.html)
    \item Tafelbild: Konjugationstabelle hablar
    \item Bildkarten: Aktivitäten (hablar, estudiar, etc.)
\end{enumerate}

% ============================================================================
% 10. DIDAKTIK-SCORE
% ============================================================================
\newpage
\section{Didaktik-Score}

\begin{scorebox}
\textbf{Bewertung der didaktischen Qualität nach 10 Dimensionen}\\
\textbf{Ziel-Score: $\geq$ 9.0 (Premium-Qualität)}

\vspace{1em}
\begin{tabular}{|c|l|c|c|p{4.5cm}|}
\hline
\textbf{\#} & \textbf{Dimension} & \textbf{Gew.} & \textbf{Score} & \textbf{Notiz} \\
\hline
1 & Differenzierung & 15\% & 9/10 & [B] nur 4 Pers., [E] alle 6 + freie Texte \\
\hline
2 & Taxonomie (AFB) & 10\% & 9/10 & I: Nennen, II: Anwenden, III: Texte \\
\hline
3 & Lernziel-Alignment & 10\% & 10/10 & Rahmenplan Grammatik erfüllt \\
\hline
4 & Aktivierung & 10\% & 9/10 & Chorisch, PA-Interview, Schreiben \\
\hline
5 & Scaffolding & 10\% & 9/10 & Tabelle, Beispiele, gestufte Aufgaben \\
\hline
6 & Feedback-Qualität & 10\% & 9/10 & Mündlich + LP mit Auto-Check \\
\hline
7 & Sprachliche Klarheit & 10\% & 10/10 & Altersgerecht, kurze Sätze \\
\hline
8 & Motivation \& Relevanz & 10\% & 9/10 & Schulalltag = Lebenswelt \\
\hline
9 & Inklusion (UDL) & 10\% & 9/10 & Visuell + auditiv + kinästhetisch \\
\hline
10 & Praktikabilität & 5\% & 10/10 & 90 Min. realistisch \\
\hline
\multicolumn{3}{|r|}{\textbf{GESAMT:}} & \textbf{9.2/10} & \textcolor{successgreen}{\textbf{FREIGABE}} \\
\hline
\end{tabular}

\vspace{1em}
\textbf{Bewertungsschwellen:}
\begin{itemize}
    \item[$\boxtimes$] \textcolor{successgreen}{$\geq$ 9.0: \textbf{FREIGABE} (Premium-Qualität)}
    \item[$\Box$] \textcolor{warningorange}{8.0--8.9: Iteration empfohlen}
    \item[$\Box$] 6.0--7.9: Überarbeitung erforderlich
    \item[$\Box$] $<$ 6.0: Grundlegende Neukonzeption
\end{itemize}
\end{scorebox}

\end{document}
